problem:  
Let \( (G/H,g) \) be a connected Riemannian homogeneous space with reductive decomposition  
\(\mathfrak g = \mathfrak h \oplus \mathfrak m\), and \(g\) induced by an \( \mathrm{Ad}(H)\)–invariant inner product on \(\mathfrak m\).  

For \(X \in \mathfrak m\), let \(\gamma_X(t)=\exp(tX)\cdot o\) denote the corresponding homogeneous geodesic whenever \(X\) satisfies the geodesic lemma. Write \(R_t\) for the Riemannian curvature operator at \(\gamma_X(t)\):
\[
R_t(Y,Z)W = R(Y,Z)W|_{\gamma_X(t)},
\]
viewed as a tensor field along \(\gamma_X\). Define \(J_t := P_t^{-1}\circ (R_t(\gamma_X'(t),\cdot)\gamma_X'(t))\circ P_t:\mathfrak m\to\mathfrak m\), where \(P_t:T_oM\to T_{\gamma_X(t)}M\) is parallel transport along \(\gamma_X\).

---

**Open problem:**  
Assume that for every geodesic vector \(X\in \mathfrak m\), the spectral decomposition of the **parallel–transported Jacobi operator**
\[
J_t = P_t^{-1}\circ R_t(\gamma_X'(t),\cdot)\gamma_X'(t)\circ P_t
\]
is constant in \(t\), i.e. the eigenvalues and their multiplicities of \(J_t\) are independent of \(t\).

1. Does this spectral constancy of the Jacobi operator along all homogeneous geodesics imply that \((G/H,g)\) is a geodesic orbit space?  
2. If not, does it at least imply that the metric is naturally reductive?  
3. Are there examples of non‑naturally‑reductive, non‑g.o. homogeneous spaces exhibiting this spectral constancy property?

---

Why is it a "good" problem:  
This corrected formulation replaces the ambiguous and ill‑defined “Nomizu operator along geodesics’’ with a canonical geometric object: the **Jacobi operator**, derived from curvature and well–defined along any geodesic via parallel transport.  
The constancy of its spectrum describes a subtle rigidity: the infinitesimal behavior of nearby geodesics along every homogeneous geodesic is spectrally invariant under parallel translation.

This property sits between classical symmetric‑space curvature parallelism (\(\nabla R = 0\)) and the purely algebraic geodesic orbit condition. It brings together **Riemannian curvature evolution**, **Lie algebraic symmetries**, and **spectral geometry**, posing a compact but profound question: does spectral rigidity of curvature along homogeneous geodesics force algebraic homogeneity of geodesic flow?  

An affirmative answer would reveal a new “spectral signature’’ for the g.o. condition; a negative one would suggest exotic homogeneous geometries with hidden curvature symmetries. The statement is clean and intrinsic—expressible with only the curvature tensor and parallel transport—yet it opens a deep investigation bridging **Riemannian homogeneous geometry**, **spectral analysis**, and **differential‐algebraic rigidity**, fully embodying the elegance and depth of a good research‐level mathematical problem.